% Chapter 1, Section 3 _Linear Algebra_ Jim Hefferon
%  http://joshua.smcvt.edu/linearalgebra
%  2001-Jun-09
\section{简化阶梯形}
在掌握了高斯消元法的操作机制之后，
我们看到它可以用不止一种方式完成。
例如，从这个矩阵
\begin{equation*}
    \begin{mat}[r]
       2  &2  \\
       4  &3
    \end{mat}
\end{equation*}
出发，可以得到下面三个阶梯形矩阵中的任何一个。
\begin{equation*}
    \begin{mat}[r]
       2  &2  \\
       0  &-1
    \end{mat}
    \qquad
    \begin{mat}[r]
       1  &1  \\
       0  &-1
    \end{mat}
    \qquad
    \begin{mat}[r]
       2  &0  \\
       0  &-1
    \end{mat}
\end{equation*}
第一个由 $-2\rho_1+\rho_2$ 得到。
第二个来自先作 $(1/2)\rho_1$，再作 $-4\rho_1+\rho_2$。
第三个来自
$-2\rho_1+\rho_2$ 之后接 $2\rho_2+\rho_1$
（在第一个行倍加操作之后，矩阵已经成阶梯形
但这一步仍然是合法的行变换）。

在本章第一节中我们曾指出，
这带来了一些问题。
一个线性方程组的任意两个阶梯形版本，自由变量的
个数是否相同？
若是，
二者的自由变量是否恰好相同？
在本节中我们将
给出一种方法来判断一个线性方程组
能否由另一个经行变换得到。
这两个问题的答案都是“是”，
它们都将由此得出。








\subsection{高斯–若尔当消元}%
% Gaussian elimination coupled with back-substitution
% solves linear systems but it is not the only method possible.
下面是高斯消元法的一种推广，它有一些优点。

\begin{example} \label{exm:GJRedReadOffSol}
为求解
\begin{equation*}
  \begin{linsys}{3}
    x  &+  &y  &-  &2z  &=  &-2  \\
       &   &y  &+  &3z  &=  &7   \\
    x  &   &   &-  &z   &=  &-1
  \end{linsys}
\end{equation*}
我们可以像通常那样，先把它化简成阶梯形。
\begin{equation*}
  \grstep{-\rho_1+\rho_3}
    \begin{amat}[r]{3}
       1  &1  &-2 &-2  \\
       0  &1  &3  &7   \\
       0  &-1 &1  &1
    \end{amat}
  \grstep{\rho_2+\rho_3}
    \begin{amat}[r]{3}
       1  &1  &-2 &-2  \\
       0  &1  &3  &7   \\
       0  &0  &4  &8
    \end{amat}
\end{equation*}
我们可以继续进入第二阶段，
把首主元变成 \( 1 \)
\begin{equation*}
    \grstep{(1/4)\rho_3}
    \begin{amat}[r]{3}
       1  &1  &-2 &-2  \\
       0  &1  &3  &7   \\
       0  &0  &1  &2
    \end{amat}
\end{equation*}
然后再进入第三阶段，利用首主元
消去每列中所有其余的元素，
做法是向上倍加。
\begin{equation*}
  \grstep[2\rho_3+\rho_1]{-3\rho_3+\rho_2}
    \begin{amat}[r]{3}
       1  &1  &0  &2   \\
       0  &1  &0  &1   \\
       0  &0  &1  &2
    \end{amat}
  \grstep{-\rho_2+\rho_1}
    \begin{amat}[r]{3}
       1  &0  &0  &1   \\
       0  &1  &0  &1   \\
       0  &0  &1  &2
    \end{amat}
\end{equation*}
答案是 \( x=1 \)、\( y=1 \) 和 \( z=2 \)。
\end{example}
%<*Pivoting>
用一个元素去消去一列中其余的元素，称为关于该元素的
\definend{主元消去}\index{pivoting}。
%</Pivoting>

注意，第一阶段中的行倍加操作从左向右进行，
而第三阶段中的倍加操作则从右向左进行。

\begin{example}  \label{exm:GJRedReadOffSolTwo}
中间阶段把首主元变成 \( 1 \) 的那些操作
互不相干，所以我们可以把其中多个合并成单独的一步。
\begin{align*}
    \begin{amat}[r]{2}
       2   &1   &7   \\
       4   &-2  &6
    \end{amat}
  &\grstep{-2\rho_1+\rho_2}
  \begin{amat}[r]{2}
       2   &1   &7   \\
       0   &-4  &-8
    \end{amat}                                   \\
  &\grstep[(-1/4)\rho_2]{(1/2)\rho_1}
  \begin{amat}[r]{2}
       1   &1/2   &7/2   \\
       0   &1     &2
    \end{amat}                                    \\
  &\grstep{-(1/2)\rho_2+\rho_1}
  \begin{amat}[r]{2}
       1   &0   &5/2   \\
       0   &1   &2
    \end{amat}
\end{align*}
答案是 $x=5/2$ 和 $y=2$。
\end{example}

%<*GaussJordanReduction>
高斯消元法的这种推广称为
\definend{高斯–若尔当方法}\index{Gauss's Method!Gauss-Jordan Method}或
\definend{高斯–若尔当消元}。\index{linear equation!solution of!Gauss-Jordan}\index{Gauss's Method!Gauss-Jordan}
%</GaussJordanReduction>
% It goes past echelon form to a more refined, more specialized,
% matrix form.

\begin{definition}\label{def:RedEchForm}
%<*df:RedEchForm>
一个矩阵或线性方程组称为处于
\definend{简化阶梯形\/}\index{echelon form!reduced}\index{reduced echelon form}，
如果除了是阶梯形之外，每个首主元都是~$1$，
并且是其所在列中唯一的非零元素。
%</df:RedEchForm>
\end{definition}

\noindent
%<*CostRedEchForm>
用高斯–若尔当消元求解方程组的代价
是额外的算术运算。
好处是我们可以直接读出
解集的描述。
%</CostRedEchForm>

在任何阶梯形的方程组中，无论是简化阶梯形与否，
当方程组因存在矛盾方程而解集为空时，我们可以直接读出这一情况。
当方程组没有矛盾且每个变量都是某行的首变量时，
我们可以读出方程组的解集只有一个元素。
而且，当方程组没有矛盾且至少有一个变量是自由变量时，
我们可以读出方程组的解集有无穷多个元素。

然而，在简化阶梯形中，我们不仅能直接读出解集的
大小，还能读出它的描述。
当然，解集为空时描述它毫无困难。
\nearbyexample{exm:GJRedReadOffSol} 与~\ref{exm:GJRedReadOffSolTwo}
表明，在解集只有一个元素的情形，这个元素就在常数列中。
下一个例子演示如何读出无穷解集的参数化。

\begin{example}
\begin{multline*}
  \begin{amat}[r]{4}
     2  &6  &1  &2  &5  \\
     0  &3  &1  &4  &1  \\
     0  &3  &1  &2  &5
  \end{amat}
  \grstep{-\rho_2+\rho_3}
  \begin{amat}[r]{4}
     2  &6  &1  &2  &5  \\
     0  &3  &1  &4  &1  \\
     0  &0  &0  &-2 &4
  \end{amat}                                        \\
  \grstep[(1/3)\rho_2 \\ -(1/2)\rho_3]{(1/2)\rho_1}
  \repeatedgrstep[-\rho_3+\rho_1]{-(4/3)\rho_3+\rho_2}
  \repeatedgrstep{-3\rho_2+\rho_1}
  \begin{amat}[r]{4}
     1  &0  &-1/2  &0  &-9/2  \\
     0  &1  &1/3   &0  &3  \\
     0  &0  &0     &1  &-2
  \end{amat}
\end{multline*}
作为线性方程组，它就是
\begin{equation*}
  \begin{linsys}{4}
     x_1  &&     &-&1/2x_3  &&   &= &-9/2  \\
          &&x_2  &+&1/3x_3  &&   &= &3  \\
          &&     &&        &{}\hspace{.5em}{}&x_4 &= &-2
  \end{linsys}
\end{equation*}
因此解集的一个描述如下。
\begin{equation*}
  S=\set{\colvec{x_1 \\ x_2 \\ x_3 \\ x_4}
                        =\colvec[r]{-9/2 \\ 3 \\ 0 \\ -2}
                         +\colvec[r]{1/2 \\ -1/3 \\ 1 \\ 0}x_3
                        \suchthat x_3\in\Re}
\end{equation*}
\end{example}

可见，阶梯形并不是方程组的某种唯一最佳形式。
其他形式，比如简化阶梯形，各有其
优点和缺点。
与其把线性方程组（以及相应的矩阵）
想象成我们对其操作的对象，
总是朝着阶梯形这个目标推进，不如把它们
看成相互关联的，
即可以用行变换从一个得到另一个。
本小节余下的部分将展开这一想法。

\begin{lemma} \label{le:RowOpsRev}
%<*lm:RowOpsRev>
初等行变换都是可逆的。
%</lm:RowOpsRev>
\end{lemma}

\begin{proof}
%<*pf:RowOpsRev>
对任意矩阵 \( A \)，
对换两行的效果可以由再把它们对换回去而抵消，
用非零的 \( k \) 倍乘一行，可以由再用
$1/k$ 倍乘来撤销，
而把第 \( i \) 行的倍数加到第 \( j \) 行（$i\neq j$）
可以由从第 \( j \) 行减去第 \( i \) 行的同样倍数来撤销。
\begin{equation*}
      A
     \grstep{\rho_i\leftrightarrow\rho_j}
     \repeatedgrstep{\rho_j\leftrightarrow\rho_i}
      A
  \qquad
        A
     \grstep{k\rho_i}
     \repeatedgrstep{(1/k)\rho_i}
      A
  \qquad
        A
     \grstep{k\rho_i+\rho_j}
     \repeatedgrstep{-k\rho_i+\rho_j}
      A
\end{equation*}
%</pf:RowOpsRev>
（我们需要 $i\neq j$ 这一条件；
见 \nearbyexercise{exer:INotJMakesRowOpsRev}。）
\end{proof}

再说一次，我们正在发展、如今又得到引理支持的观点是：
“化简为”这一说法有误导性：~当
\( A\longrightarrow B \) 时，我们不应把 \( B \) 看作
在~\( A \) 之后，或比~$A$ 更简单。
相反，我们应当把这两个矩阵看作相互关联的。
下面是一幅图。
它把本节开头出现的那些矩阵以及它们的
简化阶梯形版本画成
一簇，彼此可互相化简。
\begin{center}
  \includegraphics{gr/mp/ch1.28}
\end{center}

%<*EquivMatrices>
我们说彼此可互相化简的矩阵，关于行可化简这一关系是等价的。
下一个结果将利用等价的定义
为此提供依据。\appendrefs{equivalence relations}
%</EquivMatrices>

\begin{lemma} \label{lm:ReducesToIsEqRel}
%<*lm:ReducesToIsEqRel>
在矩阵之间，“化简为”是一个等价关系。
%</lm:ReducesToIsEqRel>
\end{lemma}

\begin{proof}
%<*pf:ReducesToIsEqRel0>
我们必须验证以下条件：
(i)~自反性，即任何矩阵都化简到它自身；
(ii)~对称性，即若 \( A \) 化简为 \( B \)，
   则 \( B \) 化简为 \( A \)；
以及 (iii)~传递性，即若 \( A \) 化简为 \( B \) 且
      \( B \) 化简为 \( C \)，则 \( A \) 化简为 \( C \)。
%</pf:ReducesToIsEqRel0>

%<*pf:ReducesToIsEqRel1>
自反性是容易的；任何矩阵经零步操作即化简到它自身。

由前面的引理，这一关系是对称的\Dash 若
\( A \) 经某些行变换化简为 \( B \)，
那么 \( B \) 也可以通过将这些操作反过来而化简为 \( A \)。
%</pf:ReducesToIsEqRel1>

%<*pf:ReducesToIsEqRel2>
对于传递性，设 \( A \) 化简为 \( B \)，
且 \( B \) 化简为 \( C \)。
把 $A \rightarrow\cdots\rightarrow B$ 的化简步骤
接上 $B \rightarrow\cdots\rightarrow C$ 的化简步骤，
就得到从 \( A \) 到 \( C \) 的化简。
%</pf:ReducesToIsEqRel2>
\end{proof}

\begin{definition} \label{df:RowEquivalence}
%<*df:RowEquivalence>
经初等行变换而可互相化简的两个矩阵是\definend{行等价的}。\index{matrix!row equivalence}%
\index{row equivalence}\index{equivalence relation!row equivalence}
%</df:RowEquivalence>
\end{definition}

%<*RowEquivalanceClasses>
下面的图把全体矩阵的集合画成一个盒子。
在这个盒子中，每个矩阵位于一个类中。
两个矩阵在同一个类中当且仅当它们可互相化简。
这些类互不相交\Dash 没有矩阵属于两个不同的类。
我们就这样把矩阵的集合划分成了
\definend{行等价类}。\appendrefs{partitions and class representatives}\index{partition!row equivalence classes}
%</RowEquivalanceClasses>
\begin{center}
  \includegraphics{gr/mp/ch1.27}
\end{center}
\noindent 其中的一个类就是上面画出的、本节开头那些相互关联的矩阵组成的簇
（它包含所有非奇异的 $\nbyn{2}$ 矩阵）。

下一小节将证明一个矩阵的简化阶梯形是
唯一的。
用行等价关系来重新表述，
我们将证明每个矩阵都
与一个且仅一个简化阶梯形矩阵行等价。
用划分的语言来说，我们要证明的是：~每个
等价类恰含一个且仅一个简化阶梯形矩阵。
于是每个简化阶梯形矩阵都可以作为它
所在类的代表元。

